\chapter{부도확률 모델과 기대수익 점수}

\section{LightGBM 학습: \texttt{fit\_lgb}}
\textbf{하려는 것.} 신청 정보(특성 \ft{X})로 부도 여부(\ft{y}=\ft{bad})를 맞히는 나무 모델을 학습하고, 평가 연도 대출의 부도확률 $p$를 낸다.
\begin{lstlisting}
LGB_PARAMS = dict(n_estimators=2000, learning_rate=0.05, num_leaves=63, subsample=0.8, subsample_freq=1,
                  colsample_bytree=0.6, min_child_samples=200, reg_lambda=5, cat_smooth=50, n_jobs=-1, verbose=-1)

def fit_lgb(Xt, yt, seed, extra=None):
    p = dict(models.LGB_PARAMS)
    p.update(extra or {})
    u = np.random.default_rng(seed).random(len(Xt)) < 0.15
    m = lgb.LGBMClassifier(random_state=seed, **p)
    m.fit(Xt[~u], yt[~u], eval_set=[(Xt[u], yt[u])], callbacks=[lgb.early_stopping(50, verbose=False)])
    return m
\end{lstlisting}
\begin{itemize}[leftmargin=*]
\item \textbf{1--2번 줄} 설정값(고정, 튜닝 안 함). 핵심만 보면 이렇다.
  \begin{itemize}
  \item \ft{n_estimators=2000}: 나무를 최대 2000그루까지 쌓는다.
  \item \ft{learning_rate=0.05}: 나무 하나가 앞 나무들의 실수를 5\%씩만 고친다. 작을수록 천천히, 안정적으로 배운다.
  \item \ft{num_leaves=63}: 나무 하나가 손님을 최대 63개 무리로 나눈다.
  \item \ft{min_child_samples=200}: 한 무리에 최소 200건. 몇 건에 맞춘 규칙(과적합)을 막는다.
  \item \ft{subsample}, \ft{colsample_bytree}: 나무마다 행 80\%, 열 60\%만 무작위로 써서 나무들이 서로 다르게 배우게 한다.
  \end{itemize}
\item \textbf{7번 줄}: 학습 표본의 15\%를 무작위로 떼어(\ft{u}=참) 검증용으로 둔다.
\item \textbf{9번 줄}: 나머지 85\%(\texttt{\textasciitilde u})로 학습한다. 검증 15\%의 성적이 50그루 연속 나아지지 않으면 멈춘다(조기 종료). 2000은 상한일 뿐이고 실제로는 수백 그루에서 멈춘다.
\end{itemize}

\section{연도마다 학습하고 $p$, 점수 만들기: 30번의 앞부분}
\begin{lstlisting}
EV = []
for Y in range(2013, 2017):
    tr = asof.training_set(d, f'{Y}-01-01')
    tri, ytr = tr.index, tr['bad'].to_numpy()
    usable = [c for c in feats if X0.loc[tri, c].nunique(dropna=True) > 1]
    evi = d.index[(d['issue_year'] == Y) & d['resolved'] & (d['term_m'] == 36)]
    ev = returns.attach_treasury(d.loc[evi]).reset_index(drop=True)
    lgd = asof.loss_given_default(tr)
    p = np.mean([fit_lgb(X0.loc[tri, usable], ytr, s).predict_proba(X0.loc[evi, usable])[:, 1]
                 for s in SEEDS], axis=0)
    ev['p'] = p
    ev['s_er'] = (1 - p) * ev['r'] - p * lgd - ev['tb']
    EV.append(ev)
ev = pd.concat(EV, ignore_index=True)
\end{lstlisting}
\begin{itemize}[leftmargin=*]
\item \textbf{2번 줄}: 평가 연도 2013, 2014, 2015, 2016을 하나씩 돈다(\ft{range(2013, 2017)}은 2017을 포함하지 않는다).
\item \textbf{3번 줄}: 그해 1월 1일 기준으로 알 수 있던 대출만으로 학습 표본을 만든다(2장).
\item \textbf{5번 줄}: 학습 표본에서 값이 한 가지뿐인 열(예: 아직 생기지 않은 컬럼)은 뺀다.
\item \textbf{6번 줄} 평가 대상: 그해 발행된, 결과가 확정된, 36개월 대출.
\item \textbf{7번 줄}: 국채 금리 \ft{tb}를 붙이고 행 번호를 0부터 다시 매긴다.
\item \textbf{8번 줄} LGD: 학습 표본 부도 대출의 평균 손실률, 즉 $1-$(받은 총액/대출액)의 평균이다.
\item \textbf{9--10번 줄}: seed 0, 1, 2로 모델을 세 번 학습해 $p$를 평균 낸다. 무작위 요소(검증 분할, 행·\allowbreak{}열 추출) 때문에 생기는 흔들림을 줄인다. \ft{predict_proba}는 [부도 아닐 확률, 부도일 확률] 두 열을 주고, \ft{[:, 1]}이 두 번째 열, 즉 부도확률이다.
\item \textbf{12번 줄} \textbf{기대수익 점수}: $s=(1-p)\,r-p\cdot\mathrm{LGD}-tb$. 안 망하면 받는 이자, 빼기 망하면 잃는 원금, 빼기 국채 기회비용.
\item \textbf{14번 줄}: 4개 연도 표를 위아래로 붙여 하나의 \ft{ev}로 만든다. 이후 모든 장의 \ft{ev}가 이것이다.
\end{itemize}

\begin{workout}[따라 해 보기: \ft{snippets.py lgbm}]
가짜 데이터 5000건(부도율 약 17\%)으로 같은 설정을 seed 세 개로 학습한다.
\begin{lstlisting}[style=out]
seed 0: 158 trees, AUC 0.982
seed 1: 182 trees, AUC 0.985
seed 2: 76 trees, AUC 0.985
3-seed mean AUC 0.984, mean p 0.161 vs actual 0.166
\end{lstlisting}
\Ans{2000그루 상한이지만 조기 종료로 76--182그루에서 멈췄다. seed마다 나무 수가 다르고, 평균을 내면 조금 더 안정적이다. 가짜 데이터라 AUC가 높고, 실제 LC 데이터에서는 0.63--0.68이다.}
\end{workout}
\Trap{\ft{predict_proba(X)}를 그대로 $p$로 쓰는 것. 두 열짜리 표이므로 \ft{[:, 1]}을 빠뜨리면 모양이 맞지 않는다. 또 최신 LightGBM은 \ft{eval_set} 대신 \ft{eval_X}, \ft{eval_y}를 권한다는 경고를 낼 수 있는데, 결과는 같다.}

\section{국채 붙이기: \texttt{attach\_treasury}}
\begin{lstlisting}
def attach_treasury(d):
    m = d.merge(treasury(), on=['ym', 'term_m'], how='left', validate='many_to_one')
    m.index = d.index
    assert m['tb'].notna().all()
    return m
\end{lstlisting}
\ft{treasury()}는 FRED에서 받은 GS3(3년), GS5(5년) 월별 금리 표다. \textbf{2번 줄}은 발행월과 만기가 같은 금리를 VLOOKUP한다. \ft{validate='many_to_one'}은 ``대출 여러 건에 금리 한 줄''이 맞는지 검사한다. \textbf{4번 줄} \ft{assert}는 ``빈 금리가 하나라도 있으면 즉시 멈춰라''라는 안전장치다.
